\newcommand{\tSS}{\mbox{$t_{\rm SS}$}}
\newcommand{\tAl}{\mbox{$t_{\rm Al}$}}
\newcommand{\tMn}{\mbox{$t_{\rm Mn}$}}
\newcommand{\tHf}{\mbox{$t_{\rm Hf}$}}
\newcommand{\tPb}{\mbox{$t_{\rm Pb}$}}
\newcommand{\alratio}{\mbox{${}^{26}{\rm Al}/{}^{27}{\rm Al}$}}
\newcommand{\mnratio}{\mbox{${}^{53}{\rm Mn}/{}^{55}{\rm Mn}$}}
\newcommand{\hfratio}{\mbox{${}^{182}{\rm Hf}/{}^{180}{\rm Hf}$}}


The chronology of the early solar system rests on extinct
(\textsuperscript{26}Al--\textsuperscript{26}Mg,
\textsuperscript{53}Mn--\textsuperscript{53}Cr,
\textsuperscript{182}Hf--\textsuperscript{182}W) and extant
(\textsuperscript{207}Pb--\textsuperscript{206}Pb) chronometers.
These should agree with each other if dating the same events, such as
isotopic closure of these systems in a rapidly cooled achondrite.  If
multiple achondrites formed on different parent bodies, the
chronometers should be concordant across them, assuming isotopically
homogeneous values of ${}^{26}{\rm Al}/{}^{27}{\rm Al}$, ${}^{53}{\rm
  Mn}/{}^{55}{\rm Mn}$ and ${}^{182}{\rm Hf}/{}^{180}{\rm Hf}$.
Testing the homogeneity of these short-lived radionuclides (SLRs) is
often done by comparing a few samples, but really requires a
statistical framework considering many meteorites simultaneously.
Here we discuss two similar statistical approaches.  The first is that
of Desch et al. (2023a,b: doi:10.1016/j.icarus.2023.115607 and
doi:10.1016/j.icarus.2023.115611), which calculates the time of
formation of an achondrite after the start of the solar system ($t=0$,
which happened a time $t_{\rm SS}$ ago), in different isotopic
systems, and takes their weighted mean as the best estimate of the
formation time.  This model was applied to 38 dates across 14
multiply-dated achondrites, adjusting Pb--Pb ages for ${}^{238}{\rm
  U}/{}^{235}{\rm U}$ fractionation during mineral formation.  We
correct an error in their calculation of the goodness-of-fit parameter
and demonstrate that the adjusted formation times show a small but
statistically significant degree of overdispersion with respect to
their analytical uncertainties.  We also introduce a second approach
to statistical chronometry, which is to calculate times of formation
before the present, essentially converting dates into `Pb--Pb times',
using a single offset parameter per isotopic system, estimated by a
weighted regression with a fixed unit slope.  Applying the second
algorithm to the same data compilation yields the same conclusion of
statistically significant overdispersion.  Extending the regression to
a random-effects model allows the excess scatter to be quantified.
The overdispersion of the Pb--Pb ages relative to the extinct
chronometers is greatly reduced after the corrections of Desch et
al. (2023) for uranium isotopic fractionation, from 450--900\,kyr to
50--380\,kyr.  Whether this remaining overdispersion is astrophysical
or analytical in origin is unresolved.  One explanation for this
remaining overdispersion might be variability of SLR abundances.  If
all variations are attributed to this, it would require 0.5 to 7.5\%
(or 40 to 130\%) variations in the ${}^{26}{\rm Al}/{}^{27}{\rm Al}$
ratio, or 2 to 7.5\% (or 5 to 22\%) variations in the ${}^{53}{\rm
  Mn}/{}^{55}{\rm Mn}$ ratio, or 1.4 to 6.8\% (or 3 to 10\%)
variations in the ${}^{182}{\rm Hf}/{}^{180}{\rm Hf}$ ratio (values in
parentheses referring to the uncorrected Pb--Pb ages).  These
relatively small variations are difficult to rule out on purely
astrophysical grounds, but would not explain the persistent
overdispersion even among angrites, which derive from the same parent
body.  This suggests that systematic variations in analysis are at
least partly responsible for the overdispersion of achondrite
formation times.  Routine measurement of ${}^{238}{\rm U}/{}^{235}{\rm
  U}$ in the mineral phases used for Pb--Pb dating, and improvements
in development of isochrons, would test this interpretation.
